\documentclass[11pt,a4paper]{article}
\usepackage[utf8]{inputenc}
\usepackage[spanish,es-tabla]{babel}
\usepackage{amsmath,amsfonts,amssymb}
\usepackage{dsfont} % Para \mathds{1} (función indicadora)
\usepackage{geometry}
\usepackage{graphicx}
\usepackage{xcolor}
\usepackage{enumitem}
\usepackage{tikz}
\usetikzlibrary{positioning,shapes,arrows.meta,calc,babel}

\geometry{margin=2.5cm}

\title{\textbf{Soluci\'on Examen Final 2022\\TAF ISC\\IMT Atlantique}}
\author{}
\date{}

\begin{document}

\maketitle

\section*{Ejercicio 1 - Canal discreto equivalente}

\textit{\textbf{Contexto del ejercicio:}} Analizaremos un sistema de comunicaciones completo desde la emisión hasta la recepción, pasando por un canal multi-trayecto. El objetivo es determinar el \textbf{canal discreto equivalente} que modela todo el sistema en tiempo discreto.

\textit{\textbf{Cadena de procesamiento:}}
\[
\text{Símbolos } s_k \xrightarrow{h_e(t)} e(t) \xrightarrow{c(t)} \xrightarrow{g_r(t)} y(t) \xrightarrow{\text{muestreo}} y_n
\]

\textbf{Datos del problema:}

\begin{itemize}
    \item \textbf{Filtro de emisión:} $h_e(t) = \mathds{1}_{[0,T]}(t)$ (pulso rectangular de duración $T$)
    \item \textbf{Filtro de recepción:} $g_r(t) = h_e(-t)$ (matched filter - versión reflejada)
    \item \textbf{Canal de propagación:} $c(t) = \delta(t-\tau_1) + \frac{1}{2}\delta(t-\tau_2) + \frac{1}{3}\delta(t-\tau_3)$
    \item Con $\tau_1 = 5T$, $\tau_2 = 8T$, $\tau_3 = 9T$ (tres trayectos con retardos diferentes)
    \item \textbf{Señal en banda base:} $e(t) = \sum_k s_k h_e(t-kT)$
    \item \textbf{Símbolos:} $s_k \in \{+1+i, +1-i, -1+i, -1-i\}$
    \item \textbf{Ritmo símbolo:} $\frac{1}{T}$ símbolos/segundo
    \item \textbf{Señal recibida muestreada:} $y_n = y(t_0 + nT)$
\end{itemize}

\subsection*{Pregunta 1: Tipo de modulación}

\textbf{Enunciado:} ¿A qué tipo de modulación corresponde la señal $e(t)$?

\vspace{0.3cm}
\textbf{Solución paso a paso:}

\textbf{Paso 1: Analizar el alfabeto de símbolos}

\textit{\textbf{¿Qué información nos da el alfabeto?}} Los símbolos son:

\[
s_k \in \{+1+i, +1-i, -1+i, -1-i\}
\]

\textit{Características:}
\begin{itemize}
    \item Son números \textbf{complejos} (parte real + parte imaginaria)
    \item Hay $M = 4$ símbolos diferentes
    \item Parte real: $\pm 1$ (dos valores posibles)
    \item Parte imaginaria: $\pm 1$ (dos valores posibles)
    \item Módulo constante: $|s_k| = \sqrt{1^2 + 1^2} = \sqrt{2}$ para todos
\end{itemize}

\textbf{Paso 2: Representación en el plano complejo}

\textit{\textbf{Constelación:}} Los puntos forman un cuadrado centrado en el origen.

\begin{center}
\begin{tikzpicture}[scale=1.5]
    % Ejes
    \draw[-latex] (-2,0) -- (2,0) node[right] {Re (canal I)};
    \draw[-latex] (0,-2) -- (0,2) node[above] {Im (canal Q)};
    
    % Puntos de la constelación
    \filldraw[blue] (1,1) circle (2pt) node[above right] {$+1+i$};
    \filldraw[blue] (1,-1) circle (2pt) node[below right] {$+1-i$};
    \filldraw[blue] (-1,1) circle (2pt) node[above left] {$-1+i$};
    \filldraw[blue] (-1,-1) circle (2pt) node[below left] {$-1-i$};
    
    % Círculo de módulo constante
    \draw[dashed,gray] (0,0) circle (1.414);
    
    % Grilla auxiliar
    \draw[dotted,gray] (-1.5,-1.5) grid (1.5,1.5);
\end{tikzpicture}
\end{center}

\textbf{Paso 3: Identificación de la modulación}

\[
\boxed{\text{QPSK (Quadrature Phase Shift Keying)}}
\]

\textit{También conocida como:} 4-QAM cuadrada o 4-PSK

\textit{\textbf{Justificación:}}
\begin{itemize}
    \item \textbf{Q} (Quadrature): Dos portadoras en cuadratura (seno y coseno)
    \item \textbf{PSK} (Phase Shift Keying): Modulación de fase (módulo constante)
    \item \textbf{4 símbolos:} Cada símbolo transporta $\log_2(4) = 2$ bits
    \item Separación angular: $90°$ entre símbolos consecutivos
\end{itemize}

\textit{\textbf{Mapeo de bits típico:}}
\begin{center}
\begin{tabular}{|c|c|}
\hline
\textbf{Bits} & \textbf{Símbolo} \\
\hline
00 & $+1+i$ \\
01 & $-1+i$ \\
11 & $-1-i$ \\
10 & $+1-i$ \\
\hline
\end{tabular}
\end{center}

(Codificación Gray - mínima distancia de Hamming entre vecinos)

\subsection*{Pregunta 2: C\'alculo de $r(t) = (h_e * g_r)(t)$}

\textbf{Enunciado:} Calcular $r(t) = (h_e * g_r)(t)$.

\vspace{0.3cm}
\textbf{Soluci\'on paso a paso:}

\textbf{Paso 1: Recordar las funciones}

\textit{\textbf{Filtro de emisi\'on:}}
\[
h_e(t) = \begin{cases}
1 & \text{si } 0 \leq t \leq T \\
0 & \text{en otro caso}
\end{cases}
\]

\textit{\textbf{Filtro de recepci\'on (matched filter):}}
\[
g_r(t) = h_e(-t) = \begin{cases}
1 & \text{si } -T \leq t \leq 0 \\
0 & \text{en otro caso}
\end{cases}
\]

\textit{\textbf{¿Por qu\'e usar matched filter?}} Maximiza la relaci\'on se\~nal-ruido (SNR) en el instante de muestreo. Es la soluci\'on \'optima cuando el ruido es blanco gaussiano.

\textbf{Paso 2: Definici\'on de convoluci\'on}

\[
r(t) = (h_e * g_r)(t) = \int_{-\infty}^{+\infty} h_e(\tau) g_r(t-\tau) d\tau
\]

\textit{Nota:} El enunciado dice $g_r(t) = h_e(t)$ pero deber\'ia ser $g_r(t) = h_e(-t)$ para matched filter. Asumimos matched filter.

Si $h_e(t) = g_r(t) = \mathds{1}_{[0,T]}(t)$:

\[
r(t) = \int_{-\infty}^{+\infty} \mathds{1}_{[0,T]}(\tau) \mathds{1}_{[0,T]}(t-\tau) d\tau
\]

\textbf{Paso 3: Determinar cu\'ando hay solapamiento}

\textit{\textbf{Condiciones para que el integrando sea no nulo:}}
\begin{itemize}
    \item $h_e(\tau) \neq 0$ requiere: $0 \leq \tau \leq T$
    \item $g_r(t-\tau) \neq 0$ requiere: $0 \leq t-\tau \leq T$ \quad $\Rightarrow$ \quad $t-T \leq \tau \leq t$
\end{itemize}

\textit{\textbf{Intersecci\'on de intervalos:}}
\[
[0, T] \cap [t-T, t] = [\max(0, t-T), \min(T, t)]
\]

Esta intersecci\'on es no vac\'ia si: $\max(0, t-T) < \min(T, t)$

\textbf{Paso 4: An\'alisis por casos seg\'un $t$}

\textit{\textbf{Visualizaci\'on:}} Imagina dos ventanas rectangulares desplaz\'andose una respecto a la otra.

\textit{\textbf{Caso 1: $t < 0$}}

No hay solapamiento entre $[0,T]$ y $[t-T, t]$ porque $t < 0$.
\[
r(t) = 0 \quad \text{(las ventanas no se tocan)}
\]

\textit{\textbf{Caso 2: $0 \leq t < T$}}

La intersecci\'on es $[0, t]$ (longitud = $t$)
\[
r(t) = \int_0^t 1 \cdot 1 \, d\tau = t
\]

\textit{Interpretaci\'on:} La funci\'on crece linealmente - las ventanas empiezan a solaparse.

\textit{\textbf{Caso 3: $T \leq t < 2T$}}

La intersecci\'on es $[t-T, T]$ (longitud = $T - (t-T) = 2T - t$)
\[
r(t) = \int_{t-T}^T 1 \cdot 1 \, d\tau = T - (t-T) = 2T - t
\]

\textit{Interpretaci\'on:} La funci\'on decrece linealmente - las ventanas empiezan a separarse.

\textit{\textbf{Caso 4: $t \geq 2T$}}

No hay solapamiento porque las ventanas est\'an completamente separadas.
\[
r(t) = 0 \quad \text{(ya no se tocan)}
\]

\textbf{Paso 5: Resultado final}

\[
\boxed{r(t) = \begin{cases}
t & \text{si } 0 \leq t < T \\
2T - t & \text{si } T \leq t < 2T \\
0 & \text{en otro caso}
\end{cases}}
\]

\textit{\textbf{Forma compacta:}} $r(t) = \Lambda\left(\frac{t-T}{T}\right) \cdot T$ donde $\Lambda$ es la funci\'on triangular.

\textit{\textbf{Propiedades clave:}}
\begin{itemize}
    \item \textbf{Forma:} Triangular (no rectangular como los pulsos originales)
    \item \textbf{Duraci\'on total:} $2T$ (de $0$ a $2T$)
    \item \textbf{Pico m\'aximo:} $r(T) = T$ (en el centro)
    \item \textbf{Simetr\'ia:} $r(T+\tau) = r(T-\tau)$ (sim\'etrica respecto a $t=T$)
    \item \textbf{\'Area:} $\int_0^{2T} r(t) dt = T^2$
    \item \textbf{Ancho de banda:} M\'as estrecho que pulso rectangular (menos lóbulos laterales)
\end{itemize}

\textit{\textbf{¿Por qu\'e es importante?}} Esta es la respuesta del sistema emisor+receptor sin canal. Define las propiedades de ISI del sistema.

\subsection*{Pregunta 3: Dibujar $(r * c)(t)$}

\textbf{Enunciado:} Dibujar $(r * c)(t)$.

\vspace{0.3cm}
\textbf{Soluci\'on paso a paso:}

\textbf{Paso 1: Aplicar la propiedad de convoluci\'on con deltas}

\textit{\textbf{Propiedad fundamental:}} La convoluci\'on con una delta desplaza la funci\'on:
\[
f(t) * \delta(t-\tau_0) = f(t-\tau_0)
\]

\textit{\textbf{¿Qu\'e significa f\'isicamente?}} El canal multi-trayecto crea copias retardadas y atenuadas de la se\~nal.

Con:
\[
c(t) = \delta(t-\tau_1) + \frac{1}{2}\delta(t-\tau_2) + \frac{1}{3}\delta(t-\tau_3)
\]

Entonces:
\[
(r * c)(t) = r(t-\tau_1) + \frac{1}{2}r(t-\tau_2) + \frac{1}{3}r(t-\tau_3)
\]

\textbf{Paso 2: Sustituir los valores de retardos}

\textit{Datos:} $\tau_1 = 5T$, $\tau_2 = 8T$, $\tau_3 = 9T$

\[
(r * c)(t) = r(t-5T) + \frac{1}{2}r(t-8T) + \frac{1}{3}r(t-9T)
\]

\textit{\textbf{Interpretaci\'on:}}
\begin{itemize}
    \item \textbf{Primer trayecto:} Llega en $t = 5T$ con amplitud completa (camino m\'as fuerte)
    \item \textbf{Segundo trayecto:} Llega en $t = 8T$ con amplitud $1/2$ (reflexi\'on/dispersi\'on)
    \item \textbf{Tercer trayecto:} Llega en $t = 9T$ con amplitud $1/3$ (trayecto m\'as d\'ebil)
\end{itemize}

\textbf{Paso 3: Descripci\'on detallada de cada componente}

\textit{\textbf{Tri\'angulo 1:}} $r(t-5T)$ (Trayecto principal)
\begin{itemize}
    \item \textbf{Soporte:} $[5T, 7T]$ (duraci\'on $2T$)
    \item \textbf{Pico:} En $t = 6T$, altura $= T$
    \item \textbf{Expresi\'on:} 
    \[
    r(t-5T) = \begin{cases}
    t - 5T & \text{si } 5T \leq t < 6T \\
    7T - t & \text{si } 6T \leq t < 7T
    \end{cases}
    \]
    \item \textcolor{blue}{\textbf{Azul oscuro en el diagrama}}
\end{itemize}

\textit{\textbf{Tri\'angulo 2:}} $\frac{1}{2}r(t-8T)$ (Primer eco)
\begin{itemize}
    \item \textbf{Soporte:} $[8T, 10T]$ (duraci\'on $2T$)
    \item \textbf{Pico:} En $t = 9T$, altura $= T/2$
    \item \textbf{Expresi\'on:}
    \[
    \frac{1}{2}r(t-8T) = \begin{cases}
    \frac{t - 8T}{2} & \text{si } 8T \leq t < 9T \\
    \frac{10T - t}{2} & \text{si } 9T \leq t < 10T
    \end{cases}
    \]
    \item \textcolor{green}{\textbf{Verde en el diagrama}}
\end{itemize}

\textit{\textbf{Tri\'angulo 3:}} $\frac{1}{3}r(t-9T)$ (Segundo eco)
\begin{itemize}
    \item \textbf{Soporte:} $[9T, 11T]$ (duraci\'on $2T$)
    \item \textbf{Pico:} En $t = 10T$, altura $= T/3$
    \item \textbf{Expresi\'on:}
    \[
    \frac{1}{3}r(t-9T) = \begin{cases}
    \frac{t - 9T}{3} & \text{si } 9T \leq t < 10T \\
    \frac{11T - t}{3} & \text{si } 10T \leq t < 11T
    \end{cases}
    \]
    \item \textcolor{red}{\textbf{Rojo en el diagrama}}
\end{itemize}

\textbf{Paso 4: An\'alisis de solapamiento}

\textit{\textbf{¡Observaci\'on cr\'itica!}} Los tri\'angulos 2 y 3 se solapan en el intervalo $[9T, 10T]$:

\textit{En el intervalo $9T \leq t < 10T$:}
\[
(r*c)(t) = \underbrace{0}_{\text{tri. 1}} + \underbrace{\frac{10T-t}{2}}_{\text{tri. 2 bajando}} + \underbrace{\frac{t-9T}{3}}_{\text{tri. 3 subiendo}}
\]

\textit{Valor en $t = 9T$:} $0 + \frac{T}{2} + 0 = \frac{T}{2}$

\textit{Valor en $t = 10T$:} $0 + 0 + \frac{T}{3} = \frac{T}{3}$

\textit{\textbf{¿Hay un m\'aximo local?}} Derivando:
\[
\frac{d}{dt}\left[\frac{10T-t}{2} + \frac{t-9T}{3}\right] = -\frac{1}{2} + \frac{1}{3} = -\frac{1}{6} < 0
\]

La suma decrece mon\'otonamente - no hay pico adicional.

\textbf{Paso 5: Caracter\'isticas importantes del resultado}

\begin{itemize}
    \item \textbf{Duraci\'on total:} De $5T$ a $11T$ $\to$ extensi\'on de $6T$
    \item \textbf{Pico principal:} En $t = 6T$ con amplitud $T$
    \item \textbf{Retardo m\'inimo:} $\tau_{min} = 5T$ (primer trayecto)
    \item \textbf{Retardo m\'aximo:} $\tau_{max} = 11T$ (fin del tercer trayecto)
    \item \textbf{Delay spread:} $\Delta\tau = \tau_{max} - \tau_{min} = 6T$
    \item \textbf{Zonas de solapamiento:}
    \begin{itemize}
        \item $[8T, 7T]$: Solo tri\'angulo 1 (sin solapamiento con $[8T, 10T]$)
        \item $[9T, 10T]$: Solapamiento de tri\'angulos 2 y 3
        \item $[10T, 11T]$: Solo tri\'angulo 3
    \end{itemize}
\end{itemize}

\textit{\textbf{Implicaci\'on para ISI:}} La extensi\'on de $6T$ significa que un s\'imbolo puede interferir con hasta 6 s\'imbolos posteriores si transmitimos a ritmo $1/T$.

\subsection*{Pregunta 4: Expresi\'on de $y_n$ con ISI}

\textbf{Enunciado:} Dar la expresi\'on de $y_n = y(t_0 + nT)$ haciendo aparecer el t\'ermino de interferencia entre s\'imbolos.

\vspace{0.3cm}
\textbf{Soluci\'on paso a paso:}

\textbf{Paso 1: Expresi\'on general de la se\~nal recibida}

\textit{\textbf{Se\~nal continua recibida:}}
\[
y(t) = \sum_k s_k (r * c)(t - kT) + \text{ruido}(t)
\]

\textit{\textbf{¿Qu\'e representa cada t\'ermino?}}
\begin{itemize}
    \item $s_k$: S\'imbolo transmitido en el instante $k$
    \item $(r*c)(t-kT)$: Respuesta del sistema al s\'imbolo $s_k$
    \item La suma: superposici\'on de todas las respuestas (ISI incluida)
\end{itemize}

\textbf{Paso 2: Muestrear en $t = t_0 + nT$}

\[
y_n = y(t_0 + nT) = \sum_k s_k (r * c)(t_0 + nT - kT)
\]

Cambio de variable: $l = n-k$ $\Rightarrow$ $k = n-l$

\[
y_n = \sum_l s_{n-l} (r * c)(t_0 + lT)
\]

\textbf{Paso 3: Definir el canal discreto equivalente}

\textit{\textbf{Concepto clave - Canal discreto equivalente:}} Representa todo el sistema continuo (filtros + canal) en tiempo discreto.

Definimos los coeficientes:
\[
h_l = (r * c)(t_0 + lT)
\]

\textit{\textbf{¿Qu\'e significa $h_l$?}} Es el impacto del s\'imbolo $s_{n-l}$ sobre la muestra $y_n$.

Entonces:
\[
y_n = \sum_l s_{n-l} h_l + b_n
\]

donde $b_n$ es el ruido muestreado.

\textbf{Paso 4: Determinar el rango de $l$ (soporte del canal)}

\textit{\textbf{Pregunta:}} ¿Cu\'antos coeficientes $h_l$ son no nulos?

Recordamos que $(r * c)(t) \neq 0$ solo en $[5T, 11T]$.

Para que $h_l = (r*c)(t_0 + lT) \neq 0$, necesitamos:
\[
5T \leq t_0 + lT \leq 11T
\]

\textit{Si elegimos $t_0 = 6T$ (instante \'optimo):}
\[
5T \leq 6T + lT \leq 11T \quad \Rightarrow \quad -1 \leq l \leq 5
\]

Pero verificando los cálculos: $h_{-1} = h_1 = h_2 = h_5 = 0$

Por tanto: Solo $l \in \{0, 3, 4\}$ tienen valores no nulos $\to$ \textbf{3 coeficientes no nulos}

\textbf{Paso 5: Separar s\'imbolo deseado e ISI}

\textit{\textbf{Descomposici\'on fundamental:}}
\[
y_n = \underbrace{s_n h_0}_{\text{s\'imbolo deseado}} + \underbrace{\sum_{\substack{l \in \{3,4\}}} s_{n-l} h_l}_{\text{Interferencia ISI}} + b_n
\]

Simplificando (solo coeficientes no nulos):
\[
y_n = T \cdot s_n + \frac{T}{2} s_{n-3} + \frac{T}{3} s_{n-4} + b_n
\]

\textit{\textbf{Clasificaci\'on de la ISI:}}
\begin{itemize}
    \item \textbf{ISI post-cursor} ($l < 0$): No hay (todos los $h_l$ con $l<0$ son cero)
    \item \textbf{ISI pre-cursor} ($l > 0$): Solo contribuyen $l=3$ y $l=4$
    \[
    \text{Pre-ISI} = s_{n-3} h_3 + s_{n-4} h_4 = \frac{T}{2} s_{n-3} + \frac{T}{3} s_{n-4}
    \]
\end{itemize}

\textit{\textbf{¿Por qu\'e es importante distinguirlas?}}
\begin{itemize}
    \item Post-ISI: Requiere t\'ecnicas no causales o retardo adicional para cancelar
    \item Pre-ISI: Puede cancelarse con DFE (Decision Feedback Equalizer) - m\'as f\'acil
\end{itemize}

\textbf{Respuesta final:}

\[
\boxed{y_n = T \cdot s_n + \frac{T}{2} s_{n-3} + \frac{T}{3} s_{n-4} + b_n}
\]

Con \textbf{2 t\'erminos de interferencia ISI} (pre-cursor) adem\'as del s\'imbolo deseado.

\subsection*{Pregunta 5: Valor \'optimo de $t_0$}

\textbf{Enunciado:} ¿Qué valor de $t_0$ propone?

\vspace{0.3cm}
\textbf{Solución paso a paso:}

\textbf{Paso 1: Criterio de optimización}

\textit{\textbf{Objetivo:}} Maximizar $h_0$ (coeficiente del símbolo deseado) para mejorar la relación señal-interferencia (SIR).

\textit{\textbf{Estrategia:}} Muestrear donde $(r*c)(t)$ alcanza su máximo global.

\textbf{Paso 2: Analizar $(r * c)(t)$}

Recordamos que:
\[
(r * c)(t) = r(t-5T) + \frac{1}{2}r(t-8T) + \frac{1}{3}r(t-9T)
\]

\textit{Máximos individuales de cada componente:}
\begin{itemize}
    \item $r(t-5T)$ tiene máximo $T$ en $t = 6T$ (trayecto principal)
    \item $r(t-8T)$ tiene máximo $\frac{T}{2}$ en $t = 9T$ (primer eco)
    \item $r(t-9T)$ tiene máximo $\frac{T}{3}$ en $t = 10T$ (segundo eco)
\end{itemize}

\textbf{Paso 3: Evaluaciones comparativas}

\textit{\textbf{Candidato 1: $t_0 = 6T$}}
\begin{align*}
h_0 &= (r*c)(6T) = r(T) + \frac{1}{2}r(-2T) + \frac{1}{3}r(-3T) \\
&= T + 0 + 0 = T
\end{align*}

\textit{\textbf{Candidato 2: $t_0 = 9T$}}
\[
h_0 = (r*c)(9T) = r(4T) + \frac{1}{2}r(T) + \frac{1}{3}r(0) = 0 + \frac{T}{2} + 0 = \frac{T}{2}
\]

\textit{\textbf{Candidato 3: $t_0 = 10T$}}
\[
h_0 = (r*c)(10T) = r(5T) + \frac{1}{2}r(2T) + \frac{1}{3}r(T) = 0 + 0 + \frac{T}{3} = \frac{T}{3}
\]

\textbf{Paso 4: Comparación}

\begin{center}
\begin{tabular}{|c|c|c|}
\hline
\textbf{$t_0$} & \textbf{$h_0$} & \textbf{Resultado} \\
\hline
$6T$ & $T$ & \textcolor{green}{\textbf{MÁXIMO}} \\
$9T$ & $T/2$ & Submáximo \\
$10T$ & $T/3$ & Mínimo \\
\hline
\end{tabular}
\end{center}

\textbf{Conclusión:}

El instante óptimo de muestreo es:

\[
\boxed{t_0 = 6T}
\]

\textit{\textbf{Justificación:}}
\begin{itemize}
    \item Maximiza el coeficiente principal: $h_0 = T$
    \item Alinea con el pico del trayecto más fuerte
    \item Optimiza la relación señal-interferencia
\end{itemize}

\subsection*{Pregunta 6: Coeficientes del canal discreto equivalente}

\textbf{Enunciado:} Dar el valor de los coeficientes del canal discreto equivalente.

\vspace{0.3cm}
\textbf{Soluci\'on paso a paso:}

\textbf{Paso 1: Definici\'on}

Los coeficientes son $h_l = (r * c)(t_0 + lT)$ con $t_0 = 6T$.

\textbf{Paso 2: C\'alculo de cada coeficiente}

$h_l = r(6T + lT - 5T) + \frac{1}{2}r(6T + lT - 8T) + \frac{1}{3}r(6T + lT - 9T)$

$h_l = r(T + lT) + \frac{1}{2}r(-2T + lT) + \frac{1}{3}r(-3T + lT)$

\textbf{Para $l = -1$:}
\[
h_{-1} = r(T + (-1)T) + \frac{1}{2}r(-2T + (-1)T) + \frac{1}{3}r(-3T + (-1)T)
\]
\[
h_{-1} = r(0) + \frac{1}{2}r(-3T) + \frac{1}{3}r(-4T) = 0 + 0 + 0 = 0
\]

\textit{Nota:} $r(0) = 0$ porque el triángulo $r(t)$ comienza en $t=0$ (valor en el borde).

\textbf{Para $l = 0$:}
\[
h_0 = r(T) + \frac{1}{2}r(-2T) + \frac{1}{3}r(-3T) = T + 0 + 0 = T
\]

\textbf{Para $l = 1$:}
\[
h_1 = r(2T) + \frac{1}{2}r(-T) + \frac{1}{3}r(-2T) = 0 + 0 + 0 = 0
\]

\textbf{Para $l = 2$:}
\[
h_2 = r(3T) + \frac{1}{2}r(0) + \frac{1}{3}r(-T) = 0 + 0 + 0 = 0
\]

\textbf{Para $l = 3$:}
\[
h_3 = r(T + 3T) + \frac{1}{2}r(-2T + 3T) + \frac{1}{3}r(-3T + 3T)
\]
\[
h_3 = r(4T) + \frac{1}{2}r(T) + \frac{1}{3}r(0) = 0 + \frac{T}{2} + 0 = \frac{T}{2}
\]
\textit{Nota:} $r(T) = T$ (pico del triángulo), pero aquí multiplicamos por $\frac{1}{2}$.

\textbf{Para $l = 4$:}
\[
h_4 = r(T + 4T) + \frac{1}{2}r(-2T + 4T) + \frac{1}{3}r(-3T + 4T)
\]
\[
h_4 = r(5T) + \frac{1}{2}r(2T) + \frac{1}{3}r(T) = 0 + 0 + \frac{T}{3} = \frac{T}{3}
\]
\textit{Nota:} Solo el tercer término contribuye: $\frac{1}{3} \cdot T = \frac{T}{3}$.

\textbf{Para $l = 5$:}
\[
h_5 = r(6T) + \frac{1}{2}r(3T) + \frac{1}{3}r(2T) = 0 + 0 + 0 = 0
\]

\textbf{Respuesta:}

Canal discreto equivalente:

\[
\boxed{H(z) = T + \frac{T}{2}z^{-3} + \frac{T}{3}z^{-4}}
\]

Coeficientes no nulos: $h_0 = T$, $h_3 = \frac{T}{2}$, $h_4 = \frac{T}{3}$

\subsection*{Pregunta 7: Delay spread del canal discreto}

\textbf{Enunciado:} ¿Cu\'al es el delay spread del canal discreto equivalente?

\vspace{0.3cm}
\textbf{Soluci\'on:}

\textbf{Definici\'on:}

El delay spread es la diferencia entre el \'ultimo y el primer coeficiente no nulo del canal.

\textbf{C\'alculo:}

\begin{itemize}
    \item Primer coeficiente no nulo: $l = 0$ con $h_0 = T$
    \item Segundo coeficiente no nulo: $l = 3$ con $h_3 = \frac{T}{2}$
    \item \'Ultimo coeficiente no nulo: $l = 4$ con $h_4 = \frac{T}{3}$
\end{itemize}

\textbf{Delay spread:}
\[
L = 4 - 0 = 4
\]

\textbf{Respuesta:}

\[
\boxed{L = 4 \text{ per\'iodos de s\'imbolo}}
\]

\textit{Nota importante:} Aunque $L=4$, solo 3 coeficientes son no nulos ($h_0, h_3, h_4$). Los coeficientes $h_1$ y $h_2$ son cero debido al espaciamiento de los retardos del canal.

\subsection*{Pregunta 8: Probabilidad de error sin ecualizador}

\textbf{Enunciado:} Sin ruido, ¿cu\'al es la probabilidad de error de s\'imbolo si no se usa ecualizador?

\vspace{0.3cm}
\textbf{Soluci\'on paso a paso:}

\textbf{Paso 1: Analizar la se\~nal recibida}

\[
y_n = T \cdot s_n + \frac{T}{2} s_{n-3} + \frac{T}{3} s_{n-4}
\]

\textbf{Paso 2: Caso peor}

Para determinar si hay errores, analizamos el caso peor donde la ISI se opone al s\'imbolo deseado.

Supongamos $s_n = 1 + i$ (para simplificar tomamos la parte real):

\textbf{Caso favorable:} $s_{n-3} = s_{n-4} = 1 + i$
\[
y_n = T(1) + \frac{T}{2}(1) + \frac{T}{3}(1) = T\left(1 + \frac{1}{2} + \frac{1}{3}\right) = \frac{11T}{6}
\]

\textbf{Caso desfavorable:} $s_{n-3} = s_{n-4} = -1 - i$
\[
y_n = T(1) + \frac{T}{2}(-1) + \frac{T}{3}(-1) = T\left(1 - \frac{1}{2} - \frac{1}{3}\right) = \frac{T}{6}
\]

\textbf{Paso 3: Verificar si hay cruce de umbral}

Para QPSK, el umbral de decisi\'on est\'a en 0.

La componente real de $y_n$ var\'ia entre $\frac{T}{6}$ y $\frac{11T}{6}$.

Como $\frac{T}{6} > 0$, \textbf{siempre se decide correctamente la parte real}.

Lo mismo ocurre con la parte imaginaria por simetr\'ia.

\textbf{Conclusi\'on:}

Aunque hay ISI significativa, no hay cruce de umbrales de decisi\'on.

\textbf{Respuesta:}

\[
\boxed{P_e = 0}
\]

Sin ruido y con este canal espec\'ifico, no hay errores incluso sin ecualizador. Sin embargo, el ecualizador sigue siendo \'util para mejorar la SNR al eliminar la ISI.

\subsection*{Pregunta 9: Esquema DFE}

\textbf{Enunciado:} Conociendo el canal discreto equivalente, proponer el esquema de bloques de un ecualizador con realimentaci\'on de decisiones (DFE).

\vspace{0.3cm}
\textbf{Soluci\'on:}

\textbf{Principio del DFE:}

El ecualizador DFE tiene dos partes:
\begin{itemize}
    \item \textbf{Feedforward:} Filtra la se\~nal recibida
    \item \textbf{Feedback:} Cancela la ISI postcursor usando decisiones previas
\end{itemize}

\textbf{Esquema de bloques:}

\begin{center}
\begin{tikzpicture}[node distance=1.5cm, auto, ->]
    % Feedforward part
    \node [draw, rectangle] (ff1) {$g_{-1}$};
    \node [draw, rectangle, right of=ff1] (ff0) {$g_0$};
    \node [draw, rectangle, right of=ff0] (ff1b) {$g_1$};
    \node [coordinate, left of=ff1] (input) {};
    \node [coordinate, above of=input, node distance=0.8cm] (inputlabel) {$y_n$};
    
    % Adder
    \node [draw, circle, right of=ff1b, node distance=2cm] (add) {$+$};
    
    % Decision
    \node [draw, rectangle, right of=add, node distance=2cm] (decision) {Decisi\'on};
    \node [coordinate, right of=decision, node distance=1.5cm] (output) {};
    \node [coordinate, above of=output, node distance=0.8cm] (outputlabel) {$\hat{s}_n$};
    
    % Feedback part
    \node [draw, rectangle, below of=decision, node distance=2cm] (d1) {D};
    \node [draw, rectangle, left of=d1] (b1) {$-b_1$};
    \node [draw, rectangle, left of=b1] (d2) {D};
    \node [draw, rectangle, left of=d2] (b2) {$-b_2$};
    \node [draw, rectangle, left of=b2] (d3) {D};
    \node [draw, rectangle, left of=d3] (b3) {$-b_3$};
    
    % Connections
    \draw[->] (input) -- (ff1);
    \draw[->] (ff1) -- (ff0);
    \draw[->] (ff0) -- (ff1b);
    \draw[->] (ff1b) -- (add);
    \draw[->] (add) -- (decision);
    \draw[->] (decision) -- (output);
    
    % Feedback connections
    \draw[->] (decision) |- (d1);
    \draw[->] (d1) -- (b1);
    \draw[->] (b1) -- (d2);
    \draw[->] (d2) -- (b2);
    \draw[->] (b2) -- (d3);
    \draw[->] (d3) -- (b3);
    
    % Feedback to adder
    \draw[->] (b1) |- (add);
    \draw[->] (b2) |- (add);
    \draw[->] (b3) |- (add);
    
\end{tikzpicture}
\end{center}

\textbf{Explicaci\'on:}

\textbf{Filtro feedforward:} Coeficientes $g_{-1}, g_0, g_1$ para compensar ISI precursor

\textbf{Filtro feedback:} Coeficientes $b_1, b_2, b_3$ para cancelar ISI postcursor (correspondiente a $h_3, h_4$)

\textbf{Decisor:} Decide el s\'imbolo m\'as cercano

\textbf{Retardos D:} Almacenan decisiones previas

\textbf{Ventajas del DFE:}
\begin{itemize}
    \item No amplifica ruido en el feedback (usa decisiones, no se\~nales ruidosas)
    \item Mejor SNR que ecualizador lineal
    \item Complejidad moderada
\end{itemize}

\textbf{Inconveniente:}
\begin{itemize}
    \item Propagaci\'on de errores: una decisi\'on err\'onea contamina s\'imbolos futuros
\end{itemize}

\section*{Ejercicio 2 - Sistema OFDM}

\textbf{Contexto del problema:}

Sistema OFDM con:
\begin{itemize}
    \item Banda total: 4096 kHz
    \item N\'umero de subportadoras: $N = 64$
    \item C\'odigo corrector de rendimiento: $r_c = \frac{1}{2}$
    \item Duraci\'on prefijo c\'iclico: $T_g = \frac{T_u}{8}$
    \item Subportadoras piloto: $\frac{1}{8}$ del total
    \item Modulaci\'on: PSK-8 en cada subportadora
\end{itemize}

\subsection*{Pregunta 1: D\'ebito binario \'util}

\textbf{Enunciado:} ¿Cu\'al es el d\'ebito binario \'util (solo bits de informaci\'on)?

\vspace{0.3cm}
\textbf{Soluci\'on paso a paso:}

\textbf{Paso 1: Bits por s\'imbolo PSK-8}

PSK-8 tiene $M = 8$ s\'imbolos, por tanto:
\[
m = \log_2(8) = 3 \text{ bits por s\'imbolo}
\]

\textbf{Paso 2: Subportadoras de datos}

Total de subportadoras: $N = 64$

Subportadoras piloto: $\frac{64}{8} = 8$

Subportadoras de datos: $N_d = 64 - 8 = 56$

\textbf{Paso 3: Bits codificados por s\'imbolo OFDM}

\[
\text{Bits codificados/OFDM} = 56 \times 3 = 168 \text{ bits}
\]

\textit{Estos son los bits que se transmiten en las 56 subportadoras de datos.}

\textbf{Paso 4: Bits de informaci\'on (aplicando rendimiento del c\'odigo)}

Con rendimiento $r_c = \frac{1}{2}$, significa que por cada 2 bits transmitidos, solo 1 es informaci\'on:
\[
\text{Bits de informaci\'on/OFDM} = 168 \times \frac{1}{2} = 84 \text{ bits}
\]

\textit{Nota:} El c\'odigo corrector $r_c = 1/2$ a\~nade redundancia: de 84 bits de informaci\'on se generan 168 bits codificados.

\textbf{Paso 5: Duraci\'on de un s\'imbolo OFDM}

\[
T_{OFDM} = T_u + T_g = T_u + \frac{T_u}{8} = \frac{9T_u}{8}
\]

\textbf{Paso 6: Calcular $T_u$}

La banda total es $B = 4096$ kHz y la separaci\'on entre subportadoras es:
\[
\Delta f = \frac{B}{N} = \frac{4096 \text{ kHz}}{64} = 64 \text{ kHz}
\]

Para OFDM, la separaci\'on entre subportadoras es:
\[
\Delta f = \frac{1}{T_u}
\]

Por tanto:
\[
T_u = \frac{1}{\Delta f} = \frac{1}{64 \text{ kHz}} = 15.625 \, \mu\text{s}
\]

\textbf{Paso 7: Duraci\'on total del s\'imbolo OFDM}

\[
T_{OFDM} = \frac{9 \times 15.625}{8} = 17.578 \, \mu\text{s}
\]

\textbf{Paso 8: D\'ebito binario \'util}

\[
R_b = \frac{\text{Bits de informaci\'on}}{T_{OFDM}} = \frac{84 \text{ bits}}{17.578125 \times 10^{-6} \text{ s}}
\]

\textit{C\'alculo detallado:}
\[
R_b = \frac{84}{17.578125 \times 10^{-6}} = \frac{84}{0.000017578125} = 4{,}778{,}666.67 \text{ bits/s}
\]

\textbf{Respuesta:}

\[
\boxed{R_b \approx 4.78 \text{ Mbps}}
\]

\textit{Verificaci\'on dimensional:} bits / s = bits/s \checkmark

\subsection*{Pregunta 2: Esparcimiento temporal m\'aximo compatible}

\textbf{Enunciado:} ¿Cu\'al es el esparcimiento temporal m\'aximo del canal compatible con este sistema?

\vspace{0.3cm}
\textbf{Soluci\'on:}

\textbf{Principio:}

El prefijo c\'iclico debe ser mayor que el esparcimiento temporal del canal para evitar ISI entre s\'imbolos OFDM.

\textbf{Condici\'on:}
\[
\tau_{max} \leq T_g
\]

\textbf{C\'alculo:}

\[
T_g = \frac{T_u}{8} = \frac{15.625 \, \mu\text{s}}{8} = 1.953 \, \mu\text{s}
\]

\textbf{Respuesta:}

\[
\boxed{\tau_{max} = 1.953 \, \mu\text{s} \approx 1.95 \, \mu\text{s}}
\]

Cualquier canal con esparcimiento temporal menor a este valor ser\'a correctamente manejado por el sistema OFDM.

\subsection*{Pregunta 3: Coherencia temporal m\'inima compatible}

\textbf{Enunciado:} ¿Cu\'al es la coherencia temporal m\'inima del canal compatible con este sistema?

\vspace{0.3cm}
\textbf{Soluci\'on:}

\textbf{Principio:}

Para que el sistema funcione correctamente, el canal debe permanecer aproximadamente constante durante al menos un s\'imbolo OFDM (para poder estimarlo y ecualizarlo).

\textbf{Condici\'on m\'inima:}
\[
T_{coh} \geq T_{OFDM}
\]

\textbf{Condici\'on recomendada (para estimaci\'on con pilotos):}

Si usamos pilotos cada $T_{OFDM}$, necesitamos que el canal sea coherente durante al menos 2-3 s\'imbolos OFDM para tener una estimaci\'on robusta.

\textbf{C\'alculo:}

Coherencia temporal m\'inima absoluta:
\[
T_{coh,min} = T_{OFDM} = 17.578 \, \mu\text{s}
\]

Coherencia temporal recomendada:
\[
T_{coh,rec} \approx 2 \cdot T_{OFDM} = 35.156 \, \mu\text{s}
\]

\textbf{Respuesta:}

\[
\boxed{T_{coh,min} = 17.578 \, \mu\text{s} \approx 17.6 \, \mu\text{s}}
\]

\textbf{Interpretaci\'on:}

Si el canal var\'ia m\'as r\'apido (menor coherencia temporal), la estimaci\'on del canal basada en pilotos ser\'a inexacta y el ecualizador no funcionar\'a correctamente.

En t\'erminos de velocidad (con frecuencia portadora $f_c$):
\[
v_{max} = \frac{c \cdot T_{coh}}{2 \lambda} = \frac{c}{2 f_c T_{coh}}
\]

\section*{Ejercicio 3 - C\'odigo convolutivo: representaci\'on de bloque}

\textbf{Contexto del problema:}

C\'odigo convolutivo $(15, 17)$ en octal con:
\begin{itemize}
    \item Registros inicializados a cero: $b_1 = b_2 = b_3 = 0$
    \item Generadores en octal: $(15)_8 = (001101)_2$ y $(17)_8 = (001111)_2$
    \item Palabra de informaci\'on: $[b_3 \, b_2 \, b_1 \, b_0] = [1 \, 0 \, 1 \, 1]$
\end{itemize}

\subsection*{Pregunta 1: Palabra c\'odigo para $[1011]$}

\textbf{Enunciado:} Dar la palabra c\'odigo $m_c = [c_0^{(1)} c_0^{(2)} c_1^{(1)} c_1^{(2)} c_2^{(1)} c_2^{(2)} c_3^{(1)} c_3^{(2)}]$ correspondiente a $[b_3 \, b_2 \, b_1 \, b_0] = [1 \, 0 \, 1 \, 1]$.

\vspace{0.3cm}
\textbf{Soluci\'on paso a paso:}

\textbf{Paso 1: Interpretar los polinomios generadores}

$(15)_8 = (001101)_2 = 1 + D^2 + D^3$ para $c^{(1)}$

$(17)_8 = (001111)_2 = 1 + D + D^2 + D^3$ para $c^{(2)}$

Esto significa:
\begin{align*}
c_k^{(1)} &= b_k \oplus b_{k-2} \oplus b_{k-3} \\
c_k^{(2)} &= b_k \oplus b_{k-1} \oplus b_{k-2} \oplus b_{k-3}
\end{align*}

donde $\oplus$ es XOR (suma m\'odulo 2).

\textbf{Paso 2: Diagrama del codificador}

\begin{center}
\begin{tikzpicture}[node distance=1.5cm, auto, >=Latex]
    \node [coordinate] (input) {};
    \node [draw, rectangle, right of=input] (d1) {$D$};
    \node [draw, rectangle, right of=d1] (d2) {$D$};
    \node [draw, rectangle, right of=d2] (d3) {$D$};
    
    \node [draw, circle, above right of=d2, node distance=2cm] (xor1) {$\oplus$};
    \node [draw, circle, below right of=d2, node distance=2cm] (xor2) {$\oplus$};
    
    \node [coordinate, above of=input, node distance=0.7cm] (inlabel) {$b_k$};
    \node [coordinate, right of=xor1, node distance=1.5cm] (out1) {$c_k^{(1)}$};
    \node [coordinate, right of=xor2, node distance=1.5cm] (out2) {$c_k^{(2)}$};
    
    \draw[->] (input) -- (d1);
    \draw[->] (d1) -- node {$b_{k-1}$} (d2);
    \draw[->] (d2) -- node {$b_{k-2}$} (d3);
    \draw (d3) -- node {$b_{k-3}$} ++(1,0);
    
    % Connections for c1
    \draw[->] (input) |- (xor1);
    \draw[->] (d2) |- (xor1);
    \draw[->] (d3) |- (xor1);
    \draw[->] (xor1) -- (out1);
    
    % Connections for c2
    \draw[->] (input) |- (xor2);
    \draw[->] (d1) |- (xor2);
    \draw[->] (d2) |- (xor2);
    \draw[->] (d3) |- (xor2);
    \draw[->] (xor2) -- (out2);
\end{tikzpicture}
\end{center}

\textbf{Paso 3: Codificar bit por bit}

Entrada: $[b_3, b_2, b_1, b_0] = [1, 0, 1, 1]$ (se procesa de derecha a izquierda)

Estado inicial: $[b_{-1}, b_{-2}, b_{-3}] = [0, 0, 0]$

\textbf{Instante $k=0$:} Entra $b_0 = 1$

Estado de registros: $[b_{-1}, b_{-2}, b_{-3}] = [0, 0, 0]$

\begin{align*}
c_0^{(1)} &= b_0 \oplus b_{-2} \oplus b_{-3} = 1 \oplus 0 \oplus 0 = 1 \\
c_0^{(2)} &= b_0 \oplus b_{-1} \oplus b_{-2} \oplus b_{-3} = 1 \oplus 0 \oplus 0 \oplus 0 = 1
\end{align*}

Salida: $(c_0^{(1)}, c_0^{(2)}) = (1, 1)$

Nuevo estado: $[1, 0, 0]$

\textbf{Instante $k=1$:} Entra $b_1 = 1$

Estado de registros: $[b_0, b_{-1}, b_{-2}] = [1, 0, 0]$

\begin{align*}
c_1^{(1)} &= b_1 \oplus b_{-1} \oplus b_{-2} = 1 \oplus 0 \oplus 0 = 1 \\
c_1^{(2)} &= b_1 \oplus b_0 \oplus b_{-1} \oplus b_{-2} = 1 \oplus 1 \oplus 0 \oplus 0 = 0
\end{align*}

Salida: $(c_1^{(1)}, c_1^{(2)}) = (1, 0)$

Nuevo estado: $[1, 1, 0]$

\textbf{Instante $k=2$:} Entra $b_2 = 0$

Estado de registros: $[b_1, b_0, b_{-1}] = [1, 1, 0]$

\begin{align*}
c_2^{(1)} &= b_2 \oplus b_0 \oplus b_{-1} = 0 \oplus 1 \oplus 0 = 1 \\
c_2^{(2)} &= b_2 \oplus b_1 \oplus b_0 \oplus b_{-1} = 0 \oplus 1 \oplus 1 \oplus 0 = 0
\end{align*}

Salida: $(c_2^{(1)}, c_2^{(2)}) = (1, 0)$

Nuevo estado: $[0, 1, 1]$

\textbf{Instante $k=3$:} Entra $b_3 = 1$

Estado de registros: $[b_2, b_1, b_0] = [0, 1, 1]$

\begin{align*}
c_3^{(1)} &= b_3 \oplus b_1 \oplus b_0 = 1 \oplus 1 \oplus 1 = 1 \\
c_3^{(2)} &= b_3 \oplus b_2 \oplus b_1 \oplus b_0 = 1 \oplus 0 \oplus 1 \oplus 1 = 1
\end{align*}

Salida: $(c_3^{(1)}, c_3^{(2)}) = (1, 1)$

\textbf{Respuesta:}

\[
\boxed{m_c = [11 \, 10 \, 10 \, 11]}
\]

O en forma lineal: $m_c = [1, 1, 1, 0, 1, 0, 1, 1]$

\subsection*{Pregunta 2: Rendimiento del c\'odigo}

\textbf{Enunciado:} ¿Cu\'al es el rendimiento del c\'odigo?

\vspace{0.3cm}
\textbf{Soluci\'on:}

\textbf{Definici\'on:}

El rendimiento (rate) de un c\'odigo es:
\[
r = \frac{\text{N\'umero de bits de entrada}}{\text{N\'umero de bits de salida}}
\]

\textbf{C\'alculo:}

Por cada bit de entrada, el codificador produce 2 bits de salida (uno por cada polinomio generador).

\[
r = \frac{1}{2}
\]

\textbf{Respuesta:}

\[
\boxed{r = \frac{1}{2}}
\]

Este c\'odigo duplica el n\'umero de bits transmitidos, a\~nadiendo redundancia para correcci\'on de errores.

\subsection*{Pregunta 3: Calcular $m_c \cdot h^T$}

\textbf{Enunciado:} Calcular $m_c \cdot h^T$ con $h = [11101111]$.

\vspace{0.3cm}
\textbf{Soluci\'on paso a paso:}

\textbf{Paso 1: Expresar los vectores}

\[
m_c = [1, 1, 1, 0, 1, 0, 1, 1]
\]

\[
h^T = [1, 1, 1, 0, 1, 1, 1, 1]^T
\]

\textbf{Paso 2: Producto escalar (en $\mathbb{F}_2$)}

\begin{align*}
m_c \cdot h^T &= (1 \cdot 1) \oplus (1 \cdot 1) \oplus (1 \cdot 1) \oplus (0 \cdot 0) \\
&\quad \oplus (1 \cdot 1) \oplus (0 \cdot 1) \oplus (1 \cdot 1) \oplus (1 \cdot 1) \\
&= 1 \oplus 1 \oplus 1 \oplus 0 \oplus 1 \oplus 0 \oplus 1 \oplus 1
\end{align*}

\textit{Conteo:} Posiciones con producto = 1: posiciones 1,2,3,5,7,8 $\to$ 6 unos

\[
= \underbrace{1 \oplus 1}_{=0} \oplus \underbrace{1 \oplus 1}_{=0} \oplus \underbrace{1 \oplus 1}_{=0} = 0
\]

(Suma de 6 unos = $0$ en m\'odulo 2)

\textbf{Respuesta:}

\[
\boxed{m_c \cdot h^T = 0}
\]

\textbf{Interpretaci\'on:}

El vector $h$ es ortogonal a la palabra c\'odigo $m_c$. Esto sugiere que $h$ podr\'ia ser una fila de la matriz de paridad $H$ del c\'odigo.

\subsection*{Pregunta 4: Matriz de paridad para 5 bits de entrada}

\textbf{Enunciado:} Proponer una matriz de paridad de tama\~no $(2 \times 10)$ para aplicar al c\'odigo obtenido si la entrada es $[b_4, b_3, b_2, b_1, b_0]$.

\vspace{0.3cm}
\textbf{Soluci\'on paso a paso:}

\textbf{Paso 1: Comprender la estructura}

Con 5 bits de entrada, la salida tiene $5 \times 2 = 10$ bits.

La palabra c\'odigo ser\'a: $m_c = [c_0^{(1)}, c_0^{(2)}, c_1^{(1)}, c_1^{(2)}, c_2^{(1)}, c_2^{(2)}, c_3^{(1)}, c_3^{(2)}, c_4^{(1)}, c_4^{(2)}]$

\textbf{Paso 2: Propiedad de las palabras c\'odigo}

Para c\'odigos convolutivos, la matriz de paridad $H$ satisface:
\[
m_c \cdot H^T = 0
\]

\textbf{Paso 3: Construir $H$ usando polinomios generadores}

Para el c\'odigo $(15, 17)$, podemos construir $H$ usando:

\textbf{Fila 1:} Verificaci\'on de paridad basada en $g_2(D)$:
\[
h_1 = [1, 1, 1, 0, 1, 1, 1, 1, 0, 0]
\]

\textbf{Fila 2:} Verificaci\'on de paridad desplazada:
\[
h_2 = [0, 0, 1, 1, 1, 0, 1, 1, 1, 1]
\]

\textbf{Paso 4: Matriz de paridad propuesta}

\[
H = \begin{pmatrix}
1 & 1 & 1 & 0 & 1 & 1 & 1 & 1 & 0 & 0 \\
0 & 0 & 1 & 1 & 1 & 0 & 1 & 1 & 1 & 1
\end{pmatrix}
\]

\textbf{Verificaci\'on:}

Esta matriz verifica que cada par de bits de salida cumple las restricciones impuestas por los polinomios generadores.

\textbf{Respuesta:}

\[
\boxed{H = \begin{pmatrix}
1 & 1 & 1 & 0 & 1 & 1 & 1 & 1 & 0 & 0 \\
0 & 0 & 1 & 1 & 1 & 0 & 1 & 1 & 1 & 1
\end{pmatrix}}
\]

\textbf{Propiedad:}

Para cualquier palabra c\'odigo v\'alida $m_c$ generada por el codificador:
\[
m_c \cdot H^T = \begin{pmatrix} 0 \\ 0 \end{pmatrix}
\]

Si recibimos una palabra $r$ con errores, el s\'indrome ser\'a:
\[
s = r \cdot H^T \neq \begin{pmatrix} 0 \\ 0 \end{pmatrix}
\]

lo que indica la presencia de errores.

\vspace{1cm}

\textbf{Conclusi\'on del examen:}

Hemos resuelto completamente los 3 ejercicios del examen Final 2022:

\begin{enumerate}
    \item \textbf{Canal discreto equivalente:} An\'alisis completo de la respuesta del sistema, c\'alculo de coeficientes del canal discreto, delay spread, y dise\~no de ecualizador DFE
    
    \item \textbf{Sistema OFDM:} C\'alculo de d\'ebito binario, esparcimiento temporal compatible, y coherencia temporal m\'inima
    
    \item \textbf{C\'odigo convolutivo:} Codificaci\'on de secuencia, rendimiento del c\'odigo, verificaci\'on de ortogonalidad, y construcci\'on de matriz de paridad
\end{enumerate}

\end{document}
